% ABOUTME: Models and tuning section of the IEEE report (issue #5).
% ABOUTME: Elastic Net, Random Forest and XGBoost theory, tuning setup, CV results and the RUL cap.
\section{Models and Tuning}
\label{sec:models}

Each training row is one cycle of one engine, and the task is to regress
its RUL, capped at 125 cycles. We compare two naive baselines with three
models of increasing flexibility: Elastic Net, a regularised linear model;
Random Forest, an average of deep trees; and XGBoost, a sum of shallow
boosted trees. Each model is trained on two feature sets: the raw set
(cycle number and the 14 informative sensors, 15 features) and the
engineered set of Section~\ref{sec:features} (142 features). All models are
built by \texttt{build\_model} in \texttt{src/rul/models.py}. Random
Forest and XGBoost are seeded with 42; Elastic Net has no random component,
since its coordinate descent updates the coefficients in a fixed cyclic
order.

\subsection{Baselines}
The mean predictor returns the mean training RUL for every row; it ignores
the sensors and sets the error that any useful model must improve on.
Ordinary least squares on the raw features is the simplest model that uses
the sensors. Every tuned model is expected to beat both. Linear regression
is also fitted on the engineered features, as an unregularised reference
for Elastic Net on the same inputs.

\subsection{Elastic Net}
Elastic Net~\cite{zou2005elasticnet} is linear regression with a penalty
that mixes the L1 and L2 norms of the coefficients. In the parametrisation
used by scikit-learn, with $n$ rows, features $\mathbf{x}_i$ and targets
$y_i$, it solves
\begin{equation}
\begin{aligned}
\min_{\beta_0, \boldsymbol\beta}\ &
\frac{1}{2n}\sum_{i=1}^{n}\bigl(y_i - \beta_0 - \boldsymbol\beta^\top
\mathbf{x}_i\bigr)^2 \\
& + \alpha\rho\,\lVert\boldsymbol\beta\rVert_1
+ \frac{\alpha(1-\rho)}{2}\,\lVert\boldsymbol\beta\rVert_2^2 ,
\end{aligned}
\end{equation}
where $\alpha \ge 0$ sets the overall strength of the penalty and
$\rho \in [0, 1]$ (\texttt{l1\_ratio}) its L1 share. The L1 term can set
coefficients exactly to zero, so it selects features. On its own (the
Lasso, $\rho = 1$) it tends to keep one feature from a group of correlated
features and drop the others more or less arbitrarily, and the choice can
change between samples. The L2 term shrinks coefficients smoothly towards
zero, spreads weight across correlated features and makes the solution
stable. Mixing the two gives sparse models that remain stable when
features are strongly correlated, as the rolling statistics of one sensor
over three windows are.

The penalty depends on the scale of each coefficient, so features are
standardised first. The scaler and the regression form one scikit-learn
\texttt{Pipeline}, so in cross validation the scaler is fitted on each
fold's training engines only. Elastic Net is linear in its inputs, so it
has high bias and low variance. On raw sensors it cannot follow the
flat-then-accelerating shape of degradation; it can represent nonlinear
behaviour only through features that already encode it, such as rolling
slopes and the health index. The penalty lowers variance further at the
cost of more bias.

\subsection{Random Forest}
A regression tree splits the feature space into boxes and predicts the mean
target in each box. A deep tree has low bias but high variance: a small
change in the training data can change its splits completely. Random
Forest~\cite{breiman2001rf} trains $B$ trees, each on a bootstrap sample of
the rows (bagging), considers only a random subset of the features at each
split, and averages the predictions. If each tree has variance $\sigma^2$
and two trees have correlation $\rho$, the variance of the average is
\begin{equation}
\rho\,\sigma^2 + \frac{1-\rho}{B}\,\sigma^2 .
\end{equation}
Adding trees removes the second term, and sampling features at each split
lowers $\rho$, so the forest keeps the low bias of deep trees while
reducing their variance. Trees are trained independently, and as trees
are added the forest's generalisation error converges~\cite{breiman2001rf}:
adding trees does not increase overfitting; the error levels off, so more
trees only cost time.

\subsection{Gradient boosting and XGBoost}
Boosting builds an additive model one tree at a time, each tree correcting
the errors of the trees before it~\cite{friedman2001gbm}:
\begin{equation}
F_m(\mathbf{x}) = F_{m-1}(\mathbf{x}) + \eta\, f_m(\mathbf{x}),
\end{equation}
where $f_m$ is a shallow tree fitted to the negative gradient of the loss
at $F_{m-1}$ and $\eta$ is the learning rate. For squared error the
negative gradient is the residual $y_i - F_{m-1}(\mathbf{x}_i)$, so each
tree is fitted to what the current model still gets wrong. Where bagging
averages independent low-bias trees to cut variance, boosting adds
dependent high-bias trees to cut bias. Each added tree fits the training
data more closely, so the number of rounds, the learning rate and the tree
depth must be limited to control variance.

XGBoost~\cite{chen2016xgboost} adds regularisation to the boosting
objective. Tree $m$ minimises
\begin{equation}
\sum_{i=1}^{n} \ell\bigl(y_i, F_{m-1}(\mathbf{x}_i) + f_m(\mathbf{x}_i)\bigr)
+ \gamma T + \frac{\lambda}{2}\lVert\mathbf{w}\rVert_2^2 ,
\end{equation}
where $T$ is the number of leaves and $\mathbf{w}$ the leaf values. The
$\gamma$ term charges for each extra leaf, and the $\lambda$ term
(\texttt{reg\_lambda}) shrinks leaf values, so a leaf supported by few rows
moves the prediction less. XGBoost uses a second order Taylor expansion of
the loss to choose splits and leaf values, and can also subsample rows
(\texttt{subsample}) and features (\texttt{colsample\_bytree}) for each
tree, which adds bagging-like randomness.

\subsection{Tuning setup}
All tuning uses the five engine-grouped folds of
Section~\ref{sec:evaluation}, so every engine is validated exactly once and
no engine appears on both sides of a split. The engineered features depend
on statistics fitted to training data; they are rebuilt in every fold from
that fold's training engines and then applied to its validation engines.
Each model is tuned separately on each feature set, and the configuration
with the lowest mean validation RMSE is selected. Elastic Net is cheap, so
its full grid is searched; for Random Forest and XGBoost a fixed number of
configurations is drawn at random from the grid with seed 42.
Table~\ref{tab:search-space} lists the search spaces and the selected
values. The test set is not used.

\begin{table}[t]
\centering
\caption{Search spaces. Selected values are in parentheses, given as
raw / engineered where the two feature sets differ. Random Forest and
XGBoost sample 6 and 8 configurations from their grids; Elastic Net
searches all 12.}
\label{tab:search-space}
\footnotesize
\begin{tabular}{llp{2.9cm}}
\hline
Model & Hyperparameter & Values (selected) \\
\hline
Elastic Net & \texttt{alpha} & 0.1, 0.3, 1.0 (0.1) \\
            & \texttt{l1\_ratio} & 0.2, 0.5, 0.8, 1.0 (1.0) \\
\hline
Random      & \texttt{n\_estimators} & 100, 200 (200 / 100) \\
Forest      & \texttt{max\_depth} & 8, 12, none (12 / none) \\
            & \texttt{min\_samples\_leaf} & 1, 5, 10 (1 / 5) \\
            & \texttt{max\_features} & 0.3, 0.5, 1.0 (0.3 / 0.5) \\
\hline
XGBoost     & \texttt{n\_estimators} & 200, 400 (400) \\
            & \texttt{learning\_rate} & 0.03, 0.05, 0.1 (0.03 / 0.05) \\
            & \texttt{max\_depth} & 3, 4, 6 (4) \\
            & \texttt{subsample} & 0.8, 1.0 (1.0) \\
            & \texttt{colsample\_bytree} & 0.5, 0.8, 1.0 (0.5 / 0.8) \\
            & \texttt{reg\_lambda} & 1, 5 (5) \\
\hline
\end{tabular}
\end{table}

\subsection{Cross validation results}
Table~\ref{tab:cv-results} reports each selected configuration. RMSE and
MAE are the mean and standard deviation over the five folds, computed on
every held-out cycle against the capped target. Late RMSE restricts the
error to cycles with true RUL of at most 50, and the last column to cycles
with true RUL of at most 125, where capped and true labels agree.

These are model-selection scores. The same folds are used to choose each
configuration and to report its score, so the scores of the selected
configurations are optimistic: among several configurations, the one that
did best on these folds is partly the one that was luckiest on them.
Grouping by engine prevents leakage between engines, and a small search
averaged over five folds limits the optimism, but neither removes this
selection bias. Nested cross validation would remove it; here the
independent estimate comes instead from the test set, which is scored once
with the frozen configurations (Section~\ref{sec:evaluation}).

\begin{table*}[t]
\centering
\caption{Five fold engine cross validation on FD001 training engines. Rows
above the line are trained on RUL capped at 125 and scored against it. The
two rows below the line retrain the best configuration with and without the
cap and score both against the true RUL, so their RMSE and MAE are
comparable with each other but not with the rows above; the last column is
comparable across all rows.}
\label{tab:cv-results}
\small
\begin{tabular}{lllrrrr}
\hline
Model & Features & Target & RMSE & MAE & Late RMSE & RMSE, RUL $\le 125$ \\
\hline
Mean baseline     & --         & capped   & $41.67 \pm 0.00$ & $36.97 \pm 0.00$ & 63.56 & 43.76 \\
Linear regression & raw        & capped   & $20.12 \pm 1.54$ & $16.37 \pm 1.29$ & 21.71 & 20.79 \\
Elastic Net       & raw        & capped   & $20.12 \pm 1.54$ & $16.37 \pm 1.29$ & 21.72 & 20.78 \\
Random Forest     & raw        & capped   & $16.43 \pm 0.77$ & $11.47 \pm 0.71$ & 15.40 & 18.29 \\
XGBoost           & raw        & capped   & $16.42 \pm 0.87$ & $11.38 \pm 0.70$ & 15.23 & 18.28 \\
Linear regression & engineered & capped   & $17.07 \pm 0.78$ & $13.72 \pm 0.63$ & 16.89 & 17.65 \\
Elastic Net       & engineered & capped   & $17.75 \pm 1.06$ & $14.33 \pm 0.85$ & 17.63 & 18.15 \\
Random Forest     & engineered & capped   & $14.42 \pm 0.48$ & $9.87 \pm 0.56$  & 10.98 & 15.76 \\
XGBoost           & engineered & capped   & $\mathbf{13.99} \pm 0.46$ & $\mathbf{9.81} \pm 0.42$ & \textbf{9.88} & \textbf{15.09} \\
\hline
XGBoost           & engineered & capped   & $47.47 \pm 2.21$ & $30.38 \pm 0.77$ & 9.88  & 15.09 \\
XGBoost           & engineered & uncapped & $38.26 \pm 1.41$ & $26.16 \pm 1.65$ & 11.20 & 25.99 \\
\hline
\end{tabular}
\end{table*}

Engineered features improve every model, by 2.0 to 2.4 RMSE, well beyond
the fold standard deviations. The tree ensembles beat the linear models by
a wide margin on both feature sets. Degradation in the sensors is
nonlinear and accelerates near failure. A linear model can follow this only
to the extent that its features already encode it: the rolling statistics,
slopes and health index supply part of that structure, which is why linear
regression improves from 20.12 to 17.07 with them. Trees learn the
remaining nonlinearity directly, piece by piece. XGBoost on
engineered features has the lowest error on every metric, though Random
Forest is within one fold standard deviation on overall RMSE; XGBoost's
clearer advantage is on late-life cycles (9.88 against 10.98).

All tuned models beat the mean baseline. Against the linear regression
baseline on raw features, both tree ensembles are better on both feature
sets and Elastic Net is better on engineered features (17.75 against
20.12). On raw features Elastic Net and linear regression are tied:
Elastic Net's RMSE is 0.0006 lower and its MAE 0.0015 higher, far inside
the fold standard deviation. With 15 features and over 16,000 training
rows per fold, the linear model is limited by bias rather than variance,
so the penalty has almost nothing to remove, and at small $\alpha$ Elastic
Net reduces to least squares. On engineered features, Elastic Net is less
accurate than unregularised linear regression on the same features (17.75
against 17.07). The selected penalty is pure L1, which makes the model
sparse by setting some coefficients exactly to zero, but it does not
improve predictive accuracy here; Elastic Net's gain over the raw baseline
comes from the engineered features, not from regularisation.

\subsection{Effect of the RUL cap}
Capping assumes that degradation is not visible early in life, so a model
should not be asked to tell 300 remaining cycles from 200. To test this,
the best configuration was retrained on the uncapped RUL with the same
folds, and both versions were scored against the true RUL (the last two
rows of Table~\ref{tab:cv-results}). The answer depends on which cycles
are scored. On cycles with true RUL of at most 125, the capped model is
far more accurate (RMSE 15.09 against 25.99), and it is also better on
late-life cycles (9.88 against 11.20). Over all cycles the uncapped model
is better (38.26 against 47.47). The capped model never sees a training
target above 125, so its predictions for young engines stay close to 125
while their true RUL is several hundred cycles. Capping changes the
training target rather than bounding the predictions: it concentrates the
model on the part of life where degradation is visible, trading accuracy
on early-life cycles, where predictions have little practical use, for
accuracy close to failure, where maintenance decisions are made. We keep
the cap for all models.

\subsection{Selected model}
XGBoost on engineered features, trained on RUL capped at 125, is carried
forward to the test set with the hyperparameters in
Table~\ref{tab:search-space}. The selected configuration of every model on
both feature sets is saved in \texttt{results/best\_params.json}, together
with the sensor list and window lengths needed to rebuild the features, so
test evaluation uses exactly the configurations chosen here.
